\documentclass[11pt,a4paper]{article}

\usepackage[a4paper,margin=20mm]{geometry}
\usepackage{lmodern}
\usepackage{booktabs}
\usepackage{tabularx}
\usepackage{listings}
\usepackage{xcolor}
\usepackage[hidelinks]{hyperref}

\setlength{\parskip}{6pt plus 2pt minus 1pt}
\setlength{\parindent}{0pt}

\lstset{
  basicstyle=\ttfamily\small,
  frame=single,
  rulecolor=\color{black!50},
  breaklines=true,
  columns=fullflexible,
  keepspaces=true,
  aboveskip=6pt,
  belowskip=10pt,
  literate=
    {á}{{\'a}}1 {à}{{\`a}}1 {ã}{{\~a}}1 {â}{{\^a}}1
    {é}{{\'e}}1 {ê}{{\^e}}1 {í}{{\'i}}1
    {ó}{{\'o}}1 {õ}{{\~o}}1 {ô}{{\^o}}1
    {ú}{{\'u}}1 {ç}{{\c{c}}}1
    {Á}{{\'A}}1 {Ã}{{\~A}}1 {É}{{\'E}}1 {Í}{{\'I}}1
    {Ó}{{\'O}}1 {Ú}{{\'U}}1 {Ç}{{\c{C}}}1
}

\begin{document}

{\Large\bfseries Trabalho de Pesquisa Operacional}\\[4pt]
Nome: Arthur Henrique

\vspace{4pt}
\hrule
\vspace{10pt}

\section{Introdução}
Este trabalho apresenta a implementação e a resolução de cinco problemas clássicos de Programação Linear Inteira (PLI) em redes: Fluxo de Custo Mínimo (PFCM), Transporte (PT), Designação (PD), Caminho Mínimo (PCM) e Fluxo Máximo (PFM). O objetivo é modelar cada formulação matematicamente, executar as instâncias propostas em sala e verificar a correspondência com as soluções de referência.

\section{Metodologia}
Os modelos foram implementados em Python utilizando a biblioteca DOcplex integrada ao solver IBM ILOG CPLEX. Cada problema foi organizado em uma pasta própria (\texttt{01\_PFCM} a \texttt{05\_PFM}), contendo um arquivo de entrada genérico (\texttt{in.txt}) e um notebook (\texttt{.ipynb}) responsável pela leitura dos dados, construção das variáveis, definição da função objetivo, inserção das restrições e impressão da solução ótima.

As variáveis de decisão foram definidas como inteiras não-negativas no PFCM, PT e PFM (com fluxo mínimo igual a zero nos arcos) e como binárias no PD e no PCM. Os resultados obtidos foram comparados diretamente com as soluções apresentadas no software LINGO.

\section{Resultados}
A tabela abaixo resume os valores ótimos obtidos pelo CPLEX para cada instância em comparação com os resultados de referência do LINGO.

\begin{table}[!htbp]
  \centering
  \begin{tabularx}{\linewidth}{l X l r r}
    \toprule
    \textbf{Pasta} & \textbf{Problema} & \textbf{Variáveis} & \textbf{FO CPLEX} & \textbf{FO LINGO} \\
    \midrule
    \texttt{01\_PFCM} & Fluxo de Custo Mínimo & Inteiras & 184,00  & 184,00  \\
    \texttt{02\_PT}   & Transporte            & Inteiras & 1330,00 & 1330,00 \\
    \texttt{03\_PD}   & Designação            & Binárias & 96,00   & 96,00   \\
    \texttt{04\_PCM}  & Caminho Mínimo        & Binárias & 22,00   & 22,00   \\
    \texttt{05\_PFM}  & Fluxo Máximo          & Inteiras & 19,00   & 19,00   \\
    \bottomrule
  \end{tabularx}
\end{table}

\subsection{Fluxo de Custo Mínimo (\texttt{01\_PFCM})}
\begin{lstlisting}
Status: integer optimal solution
Função Objetivo (Custo Mínimo): 184.00

Fluxo ótimo por arco:
  x[F->P] =  2   (custo unit.: 7, cap.: INF)
  x[F->M] =  8   (custo unit.: 2, cap.: 8)
  x[G->M] =  7   (custo unit.: 5, cap.: 7)
  x[G->H] =  3   (custo unit.: 2, cap.: INF)
  x[H->M] =  4   (custo unit.: 6, cap.: INF)
  x[H->N] =  3   (custo unit.: 5, cap.: 9)
  x[H->S] =  6   (custo unit.: 4, cap.: 6)
  x[M->P] = 16   (custo unit.: 1, cap.: 17)
  x[M->Q] =  3   (custo unit.: 3, cap.: 4)
  x[M->R] =  0   (custo unit.: 4, cap.: INF)
  x[N->S] =  3   (custo unit.: 3, cap.: INF)
  x[P->Q] = 10   (custo unit.: 1, cap.: 10)
  x[Q->R] =  6   (custo unit.: 1, cap.: INF)
  x[R->N] =  0   (custo unit.: 2, cap.: INF)
\end{lstlisting}

\subsection{Problema do Transporte (\texttt{02\_PT})}
\begin{lstlisting}
Status: integer optimal solution
Função Objetivo (Custo Mínimo): 1330.00

Unidades transportadas (Origem -> Destino):
  x[1->1] =  5   (custo unit.: 14)
  x[1->2] =  0   (custo unit.: 16)
  x[1->3] = 25   (custo unit.: 13)
  x[1->4] =  0   (custo unit.: 18)
  x[2->1] = 15   (custo unit.: 8)
  x[2->2] =  1   (custo unit.: 9)
  x[2->3] =  0   (custo unit.: 10)
  x[2->4] = 34   (custo unit.: 11)
  x[3->1] =  0   (custo unit.: 18)
  x[3->2] = 27   (custo unit.: 16)
  x[3->3] =  0   (custo unit.: 21)
  x[3->4] =  0   (custo unit.: 20)
\end{lstlisting}

\subsection{Problema da Designação (\texttt{03\_PD})}
\begin{lstlisting}
Status: integer optimal solution
Função Objetivo (Custo Mínimo): 96.00

Designações (Pessoa -> Tarefa):
  x[1->A] = 0   (custo: 40)
  x[1->B] = 0   (custo: 37)
  x[1->C] = 1   (custo: 35)
  x[2->A] = 1   (custo: 36)
  x[2->B] = 0   (custo: 38)
  x[2->C] = 0   (custo: 34)
  x[3->A] = 0   (custo: 29)
  x[3->B] = 1   (custo: 25)
  x[3->C] = 0   (custo: 26)
\end{lstlisting}

\subsection{Problema do Caminho Mínimo (\texttt{04\_PCM})}
\begin{lstlisting}
Status: integer optimal solution
Função Objetivo (Custo Mínimo): 22.00

Caminho mínimo encontrado: A -> B -> D -> C -> G

Arcos escolhidos (x_ij = 1):
  x[A->B] = 1   (custo: 10)
  x[B->D] = 1   (custo: 5)
  x[C->G] = 1   (custo: 4)
  x[D->C] = 1   (custo: 3)
\end{lstlisting}

\subsection{Problema de Fluxo Máximo (\texttt{05\_PFM})}
\begin{lstlisting}
Status: integer optimal solution
Função Objetivo (Fluxo Máximo S->T): 19.00

Fluxo ótimo por arco:
  x[S->A] =  8 / 8
  x[S->D] = 11 / 18
  x[A->B] =  8 / 12
  x[A->C] =  0 / 4
  x[A->D] =  0 / 10
  x[B->T] = 13 / 20
  x[C->B] =  5 / 7
  x[C->E] =  0 / 14
  x[D->C] =  5 / 5
  x[D->E] =  6 / 10
  x[E->T] =  6 / 6
\end{lstlisting}

\section{Conclusão}
Todos os cinco modelos foram resolvidos de forma ótima pelo CPLEX, reproduzindo exatamente os valores de função objetivo e as alocações obtidas no LINGO. A leitura padronizada por arquivos de entrada e a estrutura modular em Python permitem reutilizar as implementações para novas instâncias sem alteração na formulação dos modelos.

\end{document}
